\documentclass[11pt]{article}
\usepackage[utf8]{inputenc}
\usepackage{amsmath,amssymb,amsthm}
\usepackage[margin=1in]{geometry}
\usepackage{microtype}
\usepackage{hyperref}
\usepackage{xcolor}
\usepackage{enumitem}
\usepackage{newunicodechar}
\emergencystretch=3em
\newunicodechar{ℝ}{\ensuremath{\mathbb{R}}}
\newunicodechar{λ}{\ensuremath{\lambda}}
\newunicodechar{α}{\ensuremath{\alpha}}
\newunicodechar{κ}{\ensuremath{\kappa}}
\newunicodechar{·}{\ensuremath{\cdot}}
\newunicodechar{∂}{\ensuremath{\partial}}

\newcommand{\deriv}{\operatorname{deriv}}
\newcommand{\tideref}[2][]{\href{https://therisensea.org/tides/#2/}{\ifx&#1&\texttt{#2}\else#1\fi}}
\newcommand{\experimentref}[2][]{\href{https://therisensea.org/experiments/#2/}{\ifx&#1&\texttt{#2}\else#1\fi}}

\title{Tide retrospective:\\\emph{the large-\(t\) asymptote of the cross-susceptibility}}
\author{Claude Opus 4.8 (1M ctx), with Daniel Murfet}
\date{2026-05-31}

\begin{document}
\maketitle

\begin{abstract}
The asymptotic capstone deferred by the FDT-capstone tide: the $t\to\infty$
limit of the anharmonic cross-susceptibility,
\[
  t\cdot\frac{\partial}{\partial h}\Big|_0\mathrm{Cov}_h(x,x)
    = -t^2\kappa_3 \;\longrightarrow\; \frac{\alpha}{\lambda^3}.
\]
This is the closed-form value validated empirically by HMC in the 1D
FDT-identity experiment. Proof: compose the unconditional three-point identity
$\partial_h\mathrm{Cov}_h|_0 = -t\kappa_3$ with the third-cumulant asymptotic
$t^2\kappa_3\to-\alpha/\lambda^3$ via \texttt{Tendsto.congr'}. New file
\texttt{AnharmonicFDTAsymptotic.lean}; builds clean first try, no
\texttt{sorry}/\texttt{axiom}/\texttt{native\_decide}.
\end{abstract}

\section{Setting}

The FDT-capstone tide\footnote{\tideref{2026-05-30-tide-anharmonic-fdt-capstone}.}
made the anharmonic three-point identity $\partial_h\mathrm{Cov}_h(x,x)|_0 =
-t\,\kappa_3(x,x,x)$ unconditional but left its $t\to\infty$ asymptote as an
explicit TODO --- the closed-form $\alpha/\lambda^3$ that the 1D FDT-identity
experiment had measured\footnote{\experimentref{2026-05-23-experiment-fdt-identity-1d}:
HMC at $(\lambda,\alpha,\gamma,t)=(1,0.5,1,100)$, asymptote $0.005$ matched
within $2.5\sigma$.}. The seabed already carried the third-cumulant asymptotic
$t^2\kappa_3\to-\alpha/\lambda^3$; both it and the FDT identity are stated in
terms of \texttt{Threepoint.kappa3}, so they compose with no
namespace bridge. (The companion \emph{mean} asymptote
$\partial_h\langle x\rangle|_0\to-1/\lambda$ would need a
\texttt{Threepoint}-to-\texttt{Laplace} covariance bridge and is left as a
follow-up.)

\section{The thing formalised}

\begin{quote}\itshape
For the anharmonic potential ($\lambda,\gamma>0$, $\alpha^2<3\lambda\gamma$),
$\ t\cdot\deriv\big(h\mapsto\mathrm{Cov}_h(x,x)\big)(0)\to\alpha/\lambda^3$ as
$t\to\infty$ (a \texttt{Filter.Tendsto} at \texttt{atTop}).
\end{quote}

\section{Proof strategy}

Two lines of real content. First, $-(t^2\kappa_3)\to\alpha/\lambda^3$: negate
the third-cumulant asymptotic ($t^2\kappa_3\to-\alpha/\lambda^3$) and simplify
$-(-\alpha/\lambda^3)=\alpha/\lambda^3$ with \texttt{ring}. Second, on $t>0$
(\texttt{eventually\_gt\_atTop}) the FDT identity gives
$\deriv(\mathrm{Cov}_h)(0)=-t\kappa_3$, so $t\cdot\deriv(\mathrm{Cov}_h)(0)=
-(t^2\kappa_3)$ by \texttt{ring}; \texttt{Tendsto.congr'} transfers the limit
across this eventual equality.

\section{Roadblocks and resolutions}

None --- first-try build. The two ingredients were designed to compose: the
shared \texttt{Threepoint.kappa3} phrasing means no definitional bridge, and
\texttt{Tendsto.congr'} over \texttt{eventually\_gt\_atTop} is the standard way
to discharge the ``identity only holds for $t>0$'' caveat against an
\texttt{atTop} limit.

\section{What was Mathlib, what was new}

\textbf{Mathlib.} The \texttt{Tendsto} combinators and atTop
helper\footnote{\texttt{Filter.Tendsto.neg}, \texttt{Filter.Tendsto.congr'},
\texttt{eventually\_gt\_atTop}.}, plus \texttt{HasDerivAt.deriv} and
\texttt{ring}.

\textbf{Seabed (reused).} The unconditional FDT three-point identity and the
third-cumulant asymptotic --- both from this session / the prior anharmonic
work.

\textbf{New.} The cross-susceptibility asymptote theorem. About 30 lines.

\section{Lessons}

\begin{itemize}[itemsep=3pt]
\item \textbf{A deferred ``asymptotic TODO'' is cheap once both factors exist in
  the same namespace.} The whole tide is a \texttt{Tendsto.congr'} between a
  negated existing limit and a \texttt{ring}-rewritten existing identity; the
  hard analysis was paid for upstream.
\item \textbf{\texttt{eventually\_gt\_atTop} bridges a $t>0$ identity to an
  \texttt{atTop} limit.} The pointwise identity need only hold eventually.
\end{itemize}

\section{Follow-ups}

\begin{itemize}[itemsep=3pt]
\item \textbf{Mean asymptote $\partial_h\langle x\rangle|_0\to-1/\lambda$.}
  Needs a small definitional bridge identifying the \texttt{Threepoint}
  covariance at $h=0$ with the \texttt{Laplace} one,\footnote{To consume
  \texttt{cov\_self\_anharmonic\_asymptotic}, stated for \texttt{Laplace.gibbsCov}.}
  then the same \texttt{Tendsto.congr'} shape.
\item \textbf{Promote \texttt{amgm\_t\_abs\_x} to \texttt{lean-common}.}
\end{itemize}

\end{document}
